\documentclass[runningheads]{llncs}

% ---------------------------------------------------------------
% Include basic ECCV package

% TODO REVIEW: Insert your submission number below by replacing '*****'
% TODO FINAL: Comment out the following line for the camera-ready version
\usepackage[review,year=2026,ID=*****]{eccv}
% TODO FINAL: Un-comment the following line for the camera-ready version
%\usepackage{eccv}


% ---------------------------------------------------------------
% Other packages

% Commonly used abbreviations (\eg, \ie, \etc, \cf, \etal, etc.)
\usepackage{eccvabbrv}

% Include other packages here, before hyperref.
% NOTE: eccv.sty already loads amsmath, amssymb, subcaption, cleveref, cite.
\usepackage{graphicx}
\usepackage{booktabs}
\usepackage{bm}
\usepackage{array}
\usepackage{listings}

% The "axessiblity" package can be found at: https://ctan.org/pkg/axessibility?lang=en
\usepackage[accsupp]{axessibility}  % Improves PDF readability for those with disabilities.

% Resolve thesis figures from their original location
\graphicspath{{../../thesis/images/}}


% ---------------------------------------------------------------
% Hyperref package

% TODO FINAL: Comment out the following line for the camera-ready version
%\usepackage[pagebackref,breaklinks,colorlinks,citecolor=eccvblue]{hyperref}
% TODO FINAL: Un-comment the following line for the camera-ready version
\usepackage{hyperref}

% Support for ORCID icon
\usepackage{orcidlink}


\begin{document}

% ---------------------------------------------------------------
\title{True Absence or Detection Failure: Generative Augmentation
and Expert-Calibrated Filtering for Long-Tailed Bumblebee
Classification}

\titlerunning{Expert-Calibrated Filtering for Long-Tailed Bumblebee
Classification}

% TODO FINAL: Replace with your author list (anonymized for review).
\author{Anonymous ECCV submission}
\authorrunning{Anonymous ECCV submission}
\institute{Paper ID \#*****}

\maketitle


\begin{abstract}
Urban biodiversity monitoring increasingly relies on automated image
classification, yet rare species remain difficult to detect under
long-tailed data. We study 16 Massachusetts \emph{Bombus} species with
$\sim$$60{:}1$ imbalance, where rare taxa are often confused with
morphologically similar common species, obscuring true absence from
detection failure. We evaluate a pipeline combining morphology-guided
generation, LLM-as-judge filtering, and an expert-calibrated BioCLIP
probe trained on entomologist labels, comparing five augmentation
conditions under a fixed $+200$-image-per-rare-species budget. We find
that the synthetic--real gap in BioCLIP feature space is
species-specific rather than uniform; that an LLM-as-judge with a
structured morphology rubric is interpretable but only weakly calibrated
to expert biological validity (LLM--expert strict agreement only
$56\%$); and that an expert-supervised probe trained on $150$
entomologist labels selects subsets roughly $3{\times}$ richer in
expert-validated images at the same budget. Downstream, copy-paste
yields the only statistically supported per-species F1 gain
(\emph{B.~flavidus}) under five-fold cross-validation, while a volume
ablation from $+50$ to $+300$ generated images shows no monotonic
trend; unfiltered or LLM-filtered generation can in fact reduce
rare-tier performance. Useful augmentation depends not on image volume
alone, but on expanding visual diversity without corrupting
species-diagnostic morphology.

\keywords{Long-tailed recognition \and Fine-grained classification
\and Generative augmentation \and Synthetic-data filtering \and
Biodiversity monitoring}
\end{abstract}


% ===============================================================
\section{Introduction}
\label{sec:intro}

Cities are increasingly recognised as critical refugia for
biodiversity, hosting species displaced from degraded agricultural and
natural landscapes; for urban planning and conservation, detecting
these species before development decisions is a question of automated
monitoring reliability. When an automated system \emph{fails} to
detect a rare species, is the species truly absent, or has the system
simply missed it? Nondetection cannot be read as absence unless
detection probability approaches one~\cite{mackenzie_estimating_2002},
and rare-species recall is exactly where modern classifiers are
weakest. The stakes are concrete: standardised resurveys document
range contractions of up to $87\%$ for declining \emph{Bombus} species
across 40 U.S.\ states~\cite{cameron_patterns_2011} and a $96.4\%$
relative-abundance decline in
\emph{B.~fervidus}~\cite{jacobson_decline_2018}; whether these taxa are
extinct or simply missed by automated detectors hinges on the
classifier's behaviour on the long tail.

Automated species classification on community-collected
imagery~\cite{horn_inaturalist_2018, liu_large-scale_2019} is a
fine-grained, long-tailed problem: subordinate classes share global
appearance, individuals vary by pose, sex, caste, and viewpoint, and
observation effort is heavily skewed toward common species. We study
this problem on a 16-species \emph{Bombus} dataset of
GBIF~\cite{gbif_home_2025} imagery with a $\sim$$60{:}1$ imbalance;
three rare targets have only tens of training images each, and a
standard ResNet-50~\cite{he_deep_2015} baseline collapses on those
classes even as it remains strong on the head of the distribution.
Long-tail re-weighting and re-calibration~\cite{cao_learning_2019,
cui_class-balanced_2019, kang_decoupling_2020} reshape how a model uses
existing data, but they do not add the missing visual variation that
rare-species learning requires.

Multimodal text-to-image models can synthesise visually plausible
bumblebees from a few reference photographs, suggesting an obvious
remedy. The picture is more nuanced. Synthetic augmentation can help
fine-grained categories but is domain-dependent~\cite{he_is_2023,
michaeli_advancing_2025, zhao_ltgc_2024}; synthetic-only training
trails real-data baselines~\cite{azizi_synthetic_2023}; and generic
image-quality metrics correlate weakly with task-relevant
fidelity~\cite{borji_pros_2021, jayasumana_rethinking_2024}. For
species-level classification, an image can look like a convincing
bumblebee while carrying the wrong thoracic colour, abdominal banding,
or caste-specific structure, teaching the classifier the wrong species
signal.

This paper therefore asks not whether more synthetic images help
rare-species classification, but \emph{which} synthetic images should
be trusted as training data. We compare five augmentation conditions
under a fixed $+200$-image-per-rare-species budget, isolating filter
design from image quantity (\cref{fig:pipeline}): copy-paste of real
specimens (D2), unfiltered generation (D3), an LLM-as-judge filter with
a two-stage blind-ID and five-feature morphology rubric (D4), a
BioCLIP-centroid feature-space filter
(D5)~\cite{stevens_bioclip_2024}, and an expert-calibrated logistic
probe trained on $150$ entomologist labels (D6).

\begin{figure}[t]
  \centering
  \includegraphics[width=0.86\linewidth]{Diagram_A.jpg}
  \caption{Pipeline overview. Reference photographs and structured
  prompts feed a text-to-image model; three filters of increasing
  supervision (LLM rule, BioCLIP centroid, expert-calibrated probe)
  score the pool; volume-matched $+200$-per-species selections train a
  shared ResNet-50.}
  \label{fig:pipeline}
\end{figure}

\subsubsection{Contributions.} We contribute:
\begin{itemize}
\item \textbf{A calibration study of automated synthetic-image quality
signals.} Against $150$ entomologist labels, we show that an
LLM-as-judge morphology rubric and a BioCLIP-centroid filter are only
weakly aligned with expert biological validity, locating \emph{where}
and \emph{for which species} the signals fail.
\item \textbf{An expert-calibrated selection probe.} A logistic probe on
frozen BioCLIP with per-class F1-maximising thresholds selects
volume-matched subsets roughly $3{\times}$ richer in expert-validated
images than the LLM or centroid filters, paired with structured
morphological prompting used to generate the candidate pool.
\item \textbf{A multi-protocol empirical result.} Across bootstrap,
multi-seed, five-fold CV, and volume/subtractive ablations, better
selection does \emph{not} yield reliable downstream gains: synthetic
augmentation is bounded by fidelity to species-diagnostic morphology,
not image quantity, and does not beat simple copy-paste.
\item \textbf{Practical guidance for rare-species monitoring:} prefer
copy-paste where source images exist, gate generative augmentation on
expert-calibrated filters, and track false-absence rate rather than
held-out F1 alone.
\end{itemize}


% ===============================================================
\section{Related Work}
\label{sec:related}

\subsubsection{Long-Tailed and Fine-Grained Recognition.}
Long-tailed classification spans loss
re-weighting~\cite{cui_class-balanced_2019, lin_focal_2018,
cao_learning_2019}, decoupled training~\cite{kang_decoupling_2020}, and
multi-expert methods~\cite{wang_long-tailed_2022}. These approaches
change how a fixed dataset is used but do not introduce missing visual
variation. Fine-grained visual classification relies on attention to
localised parts and feature interactions~\cite{lin_bilinear_2018,
zheng_learning_2017, wei_fine-grained_2021}, often built on
ImageNet-pretrained backbones~\cite{he_deep_2015,
kornblith_better_2019}. Self-supervised and biology-aware foundation
models~\cite{oquab_dinov2_2024, radford_learning_2021,
stevens_bioclip_2024} provide stronger general-purpose features but do
not by themselves resolve rare-class scarcity in long-tailed
biodiversity datasets such as iNaturalist~\cite{horn_inaturalist_2018}
or BioTrove~\cite{yang_biotrove_2024}.

\subsubsection{Generative Augmentation for Rare Classes.}
Copy-paste with segmented foregrounds~\cite{ghiasi_simple_2021}
preserves authentic morphology by reusing real pixels. Game-engine and
diffusion-based synthesis can extend rare-class
coverage~\cite{beery_synthetic_2020, azizi_synthetic_2023}, but
synthetic-only training still trails real data and benefits depend on
domain~\cite{he_is_2023}. Recent work conditions generation on language
or structure to better preserve discriminative
features~\cite{michaeli_advancing_2025, zhao_ltgc_2024}.
Biology-specific generators~\cite{monsefi_taxadiffusion_2025} and
synthetic-to-real curricula~\cite{liang_diffusion_2025} reduce, but do
not eliminate, the synthetic--real gap. We complement this line by
asking how to \emph{select} from a generated pool, holding both
generator and downstream classifier fixed.

\subsubsection{Synthetic-Image Evaluation.}
FID and Inception Score~\cite{salimans_improved_2016, heusel_gans_2018}
measure global distributional similarity in ImageNet feature space and
align poorly with specialised biological
tasks~\cite{borji_pros_2021, jayasumana_rethinking_2024}. Human
evaluation is the gold standard but
expensive~\cite{otani_toward_2023}. The LLM-as-judge
paradigm~\cite{zheng_judging_2023} extends to images through learned
reward models~\cite{xu_imagereward_2023} and structured-rationale
scoring~\cite{lu_llmscore_2023}, but these systems target perceptual or
prompt-alignment quality rather than expert species-level validity. We
treat LLM judgments as one signal among several and calibrate them
against entomologist labels.

\subsubsection{Bumblebee Monitoring.}
Standardised resurveys~\cite{cameron_patterns_2011,
jacobson_decline_2018, richardson_bumble_2019} document sharp declines
across northeastern North America, and citizen-science
platforms~\cite{macphail_bumble_2024} expand coverage but inherit
species-imbalance. Deep classifiers on \emph{Bombus} images either
exclude classes with too few samples~\cite{spiesman_assessing_2021} or
collapse rare taxa into coarser
categories~\cite{bjerge_deep_2024, bjerge_hierarchical_2023}. Our
setting tightens the rare-tier constraint: rather than excluding
underrepresented species, we test whether synthetic augmentation can
keep them species-resolved.


% ===============================================================
\section{Dataset and Setup}
\label{sec:setup}

\subsubsection{Data.}
We assemble a 16-species \emph{Bombus} dataset of all taxa historically
documented in Massachusetts, downloaded from
GBIF~\cite{gbif_home_2025} via the bplusplus
collector~\cite{venverloo_towards_2024}. We restrict to adult
photographic observations from $2010$--$2025$ with
\texttt{occurrenceStatus=PRESENT}. A YOLOv8 detector pretrained on GBIF
insect imagery (confidence threshold $0.35$) crops single-specimen
$640{\times}640$ regions; multi-detection and zero-detection images are
discarded. The prepared dataset contains $15{,}630$ images
($10{,}933$ train / $2{,}335$ val / $2{,}362$ test) under a
$70/15/15$ stratified split (\cref{fig:species-distribution}). The
imbalance ratio is $59.9{:}1$ (Gini $0.377$): six common species exceed
$900$ training images each, while the rare tier \emph{B.~ashtoni},
\emph{B.~sandersoni}, and \emph{B.~flavidus} contains only $22$, $40$,
and $162$ training images respectively. These three are the targets of
all augmentation conditions.

\begin{figure}[t]
  \centering
  \includegraphics[width=0.74\linewidth]{data/species_distribution.png}
  \caption{Training-set distribution across the 16 Massachusetts
  \emph{Bombus} species (imbalance ratio $59.9{:}1$). Three rare
  targets receive synthetic augmentation; validation and test splits
  are real-image only.}
  \label{fig:species-distribution}
\end{figure}

\subsubsection{Why These Three Rare Species Are Hard.}
The targets exercise different failure modes. \emph{B.~sandersoni} has
a clean two-tone yellow-anterior, black-posterior pattern; it is
visually simple but small ($8$--$11\,\mathrm{mm}$) and easily confused
with \emph{B.~vagans}. \emph{B.~ashtoni} is a cuckoo bumblebee
(subgenus \emph{Psithyrus}, lacks corbiculae) with a predominantly
black body and a diagnostic white tail; its phenotype runs opposite to
the model's yellow-and-black worker prior. \emph{B.~flavidus} carries
extensive yellow coloration with a diagnostic yellow vertex and is
visually confusable with several other yellow-banded congeners.

\subsubsection{Classifier and Protocols.}
All conditions use the same downstream classifier: ResNet-50 pretrained
on ImageNet with a $2048{\to}512{\to}16$ head, fine-tuned end-to-end
with Adam (lr $1\!\times\!10^{-4}$, batch size $8$, ReduceLROnPlateau,
early stopping on val macro-F1, max $100$ epochs). Inputs are
$640{\times}640$ with random horizontal flip, colour jitter, and
ImageNet normalisation. We use three complementary evaluation
protocols. (i)~\emph{Fixed split with bootstrap} reports macro-F1 and
per-species F1 with 10{,}000-iteration bootstrap $95\%$ CIs over the
$2{,}362$-image test set. (ii)~\emph{Multi-seed} repeats the fixed
split with five random seeds to isolate optimisation noise, enabling
per-image flip analysis. (iii)~\emph{Five-fold CV} re-partitions the
real-image dataset into five stratified folds and is the primary
protocol for aggregate and per-species claims because it expands
effective rare-tier test coverage. Synthetic and copy-paste images are
added only to the training partition; validation and test splits are
real-image-only and identical across conditions.


% ===============================================================
\section{Method}
\label{sec:method}

\Cref{fig:pipeline} summarises the pipeline. Reference photographs and
structured prompts feed a text-to-image generator; three filters of
increasing supervision rank the generated pool; volume-matched
$+200$-image-per-species selections are added to the real training set;
a shared ResNet-50 is trained from scratch for each condition.
\Cref{fig:overview} previews the visual outputs of each augmentation
condition for the three rare targets.

\begin{figure}[t]
  \centering
  \includegraphics[width=0.76\linewidth]{data/pipeline_overview.jpg}
  \caption{Augmentation conditions for the three rare targets.
  Columns: real training image, SAM cutout, copy-paste composite
  (D2), and five viewpoint-controlled generated images (lateral,
  dorsal, three-quarter anterior/posterior, frontal).}
  \label{fig:overview}
\end{figure}

\subsection{Copy-Paste Augmentation (D2)}
\label{sec:cnp}

For each rare species, we segment foreground bee specimens from
existing training images using SAM
ViT-H~\cite{kirillov_segment_2023} and composite them onto varied
flower backgrounds with random rotation, scaling, horizontal flip, and
Gaussian boundary blending. Every pixel of the bee body comes from a
real photograph, so the augmentation re-contextualises authentic
morphology without introducing new texture; this makes copy-paste a
strong fidelity baseline against which generative methods must be
judged.

\subsection{Reference-Guided Generation with Structured Prompts}
\label{sec:gen}

Generated images are produced via reference-guided image editing
through the OpenAI \texttt{gpt-image-1.5} endpoint, pairing each prompt
with two reference photographs per rare species selected to span
different viewpoints (\eg dorsal and lateral). The reference images
provide visual anchors for body form and coloration; the prompt
specifies diagnostic traits, pose, environment, and constraints.
$500$ images per rare species ($1{,}500$ total) are produced
end-to-end with no manual pre-screening before the automated filters.
Each prompt is assembled from a structured template with eight
programmatic fields: scientific name, vernacular name, caste-specific
morphology, full species description with negative constraints, view
angle (lateral, dorsal, frontal, three-quarter anterior/posterior),
wing style (folded vs.\ slightly spread), habitat and lighting, and a
proportional scale anchor relative to in-scene flowers. Four design
principles emerged through iterative development:

\begin{enumerate}
\item \emph{Negative constraints}: explicit ``not bright yellow
thorax'' style instructions counteract the worker-bumblebee prior.
\item \emph{Front-to-back tergite colour maps}: per-region colour
sequences instead of holistic descriptors, since species-level
identification depends on \emph{where} a colour appears.
\item \emph{Proportional scale anchors}: relative size against flower
petals, since absolute size cues are unreliable.
\item \emph{Macro-crop framing}: shallow depth of field with the bee
filling the frame, matching the cropped distribution produced by the
preparation pipeline.
\end{enumerate}

\subsection{LLM-as-Judge Quality Filter (D4)}
\label{sec:llm}

A single GPT-4o judge with structured Pydantic outputs scores each
generated image in two stages.

\subsubsection{Stage 1: Blind Taxonomic Identification.}
The judge sees only the image and must identify it among the 16
candidate species, return \texttt{Unknown} for an ambiguous specimen,
or \texttt{No match} for non-bumblebee content. The target label is
hidden; this tests whether the image carries enough diagnostic
information to be recognised on its own.

\subsubsection{Stage 2: Detailed Morphology Scoring.}
The target species and its expected morphology are revealed, and the
judge scores five feature groups on a $1{-}5$ rubric: legs and
appendages (correct count, proportions, corbicula presence), wing
venation and texture, head and antennae, abdominal tergite banding, and
thoracic coloration. Score anchors range from $1$ (anatomically
impossible) to $5$ (photorealistic species-faithful match). The judge
also records diagnostic completeness (genus vs.\ species), caste
fidelity, and discrete failure-mode flags (\eg~wrong colouration,
missing/extra limbs, impossible geometry).

\subsubsection{Strict Pass Rule (D4).}
An image passes the strict rule iff (i)~the blind ID equals the target
species, (ii)~diagnostic completeness equals species, and (iii)~the
mean morphology score is $\geq 4.0$. A lenient holistic rule with mean
$\geq 3.0$ at genus-level diagnosis was found to admit too many
low-quality images to be useful as a filter and is not used
downstream.

\subsection{BioCLIP-Centroid Filter (D5)}
\label{sec:centroid}

Each generated image scores its cosine similarity to the L2-normalised
mean BioCLIP embedding of real training images for the target species,
with the pass threshold set to the species-specific median synthetic
similarity. We use BioCLIP over DINOv2 ViT-L/14 because its leave-one-out
5-NN species accuracy on real images is far higher overall ($0.657$
vs.\ $0.295$) and on the rare tier ($0.125$ vs.\ $0.072$).

\subsection{Expert-Calibrated Probe (D6)}
\label{sec:probe}

The LLM judge assesses morphology in language space and the centroid
filter assesses it in feature space, but neither is supervised by an
entomologist. We therefore collect $150$ annotations from a single
entomologist under a domain-grounded rubric;
these provide a reference standard rather than consensus ground truth
(we discuss this limitation in \cref{sec:discussion}). From them we
learn a filtering signal directly: a logistic probe on frozen
biology-aware embeddings trained on the expert pass labels, with
per-class F1-maximising thresholds from LOOCV predictions.

\subsubsection{Stratified Expert Annotation.}
Fifty generated images per rare species ($150$ total) are sampled from
four LLM-defined tiers (strict-pass, borderline, soft-fail, hard-fail)
under a floor-then-proportional rule. The entomologist completes a
two-stage interface mirroring the LLM rubric:
blind species identification at full resolution, then per-feature
morphology scoring with reference images visible. Each image yields a
strict expert label (target match, species-level diagnosis, mean morph
$\geq 4.0$) and a lenient expert label (no structural failure, genus or
species diagnosis, mean morph $\geq 3.0$).

\subsubsection{Logistic Probe on Frozen BioCLIP.}
For feature vector $z_i = \hat{\phi}(x_i)$ and binary strict expert
label $y_i\in\{0,1\}$, we fit a class-balanced logistic regression
\begin{equation}
\begin{aligned}
  \hat{p}_i &= \sigma(w^\top z_i + b),\\
  \mathcal{L}(w,b)
    &= -\!\sum_i \alpha_{y_i}\!\big[\,y_i\log\hat{p}_i
       + (1{-}y_i)\log(1{-}\hat{p}_i)\big] + \lambda\|w\|_2^2,
\end{aligned}
\end{equation}
selecting $C{=}1/\lambda$ by nested stratified CV over $\{0.001, 0.01,
0.1, 1, 10\}$ and reporting LOOCV AUC. We learn a species-specific pass
threshold $\tau_s^{\mathrm{exp}}$ that maximises F1 against the strict
label on species-restricted LOOCV predictions. A feature-configuration
ablation (BioCLIP only, LLM-only, BioCLIP$+$LLM,
BioCLIP$+$\allowbreak LLM$+$\allowbreak species one-hot) is reported in
\cref{sec:expert-results}; the deployed probe uses BioCLIP only.

\subsection{Volume Parity}
\label{sec:parity}

All three generated-image filters (D4, D5, D6) score the same
$1{,}500$-image pool ($500$ per rare species); only the selection rule
differs. After thresholding, the top-$200$ by score are retained per
species; if fewer than $200$ pass a filter's threshold, all passing
images are used (in our runs every filter exceeded the $200$ cap on
every species). All conditions, D2 through D6, therefore add exactly
$+200$ images per rare species to the real training set, decoupling
filter choice from synthetic quantity. The same generated pool and the
same filter selections are reused across the fixed split, multi-seed,
and five-fold protocols; the centroid (D5) and probe (D6) are fitted
once on the full real training set and reused across folds.


% ===============================================================
\section{Experiments}
\label{sec:experiments}

Rare-tier test support is small by construction ($n{=}6$, $10$, $36$ for
the three targets), so aggregate downstream tests have limited power and
a single prediction flip can move a per-species F1 substantially. We
therefore treat the better-powered intermediate signals---feature-space
gap, calibration against $150$ expert labels, and selection quality---as
primary evidence and read the downstream comparison conservatively.

\subsection{Baseline and Rare-Species Confusion}
\label{sec:baseline}

\Cref{tab:baseline} reports the ResNet-50 baseline on the fixed split.
Aggregate accuracy and macro-F1 ($88.2\%$, $0.815$) are dominated by
head species, while the three rare targets fall far below ($0.500$,
$0.588$, $0.623$ F1) with wide bootstrap CIs from small test support
($n{=}6$, $10$, $36$). Per-species F1 tracks training size---every
species with $\geq 200$ training images exceeds $0.75$ F1 and every one
below $200$ falls under $0.65$---and the rare-tier deficit is
concentrated in recall ($\leq 0.53$): the baseline is conservative,
missing real rare specimens rather than over-predicting.

\begin{table}[t]
  \caption{Baseline ResNet-50 on the fixed split, with
  10{,}000-iteration bootstrap $95\%$ CIs.}
  \label{tab:baseline}
  \centering
  \begin{tabular}{@{}lrrrrr@{}}
    \toprule
    Species & Tr.\,$n$ & Te.\,$n$ & P & R & F1 \\
    \midrule
    \emph{B.~ashtoni}    & 22  & 6  & 0.500 & 0.500 & 0.500 \\
    \emph{B.~sandersoni} & 40  & 10 & 0.714 & 0.500 & 0.588 \\
    \emph{B.~flavidus}   & 162 & 36 & 0.760 & 0.528 & 0.623 \\
    \midrule
    Macro avg.   & --- & 2{,}362 & --- & --- & 0.815 \\
    Accuracy     & --- & 2{,}362 & --- & --- & 0.882 \\
    \bottomrule
  \end{tabular}
\end{table}

The row-normalised confusion matrix shows three distinct error
structures that anchor the augmentation analysis: \emph{B.~sandersoni}
sends $80\%$ of its errors to a single confuser (\emph{B.~vagans},
sharing its yellow-anterior pattern), \emph{B.~flavidus} errors scatter
diffusely across $\sim$7 yellow-banded congeners, and
\emph{B.~ashtoni}'s three errors (of six test images) each go to a
different species, leaving its confusion structure statistically
unidentifiable.

\subsection{Synthetic--Real BioCLIP Gap}
\label{sec:gap}

\begin{figure}[t]
  \centering
  \includegraphics[width=0.72\linewidth]{embedding/embeddings.png}
  \caption{BioCLIP t-SNE projections. Left: real training images, where
  common species form compact clusters and rare species overlap their
  confusers. Right: real versus synthetic for the three rare targets,
  each synthetic cluster displaced from its own real cluster by a
  different amount.}
  \label{fig:tsne}
\end{figure}

\Cref{fig:tsne} projects real and synthetic images into a shared
BioCLIP space: common species form compact clusters while rare species
overlap with their confusers, matching the confusion structure above.
The median synthetic-to-centroid cosine distance is $0.31$
(\emph{ashtoni}), $0.25$ (\emph{sandersoni}), and $0.32$
(\emph{flavidus}) versus $0.10{-}0.20$ for real images---not a global
synthetic-vs-real shift but a per-species one, with visually plausible
synthetics populating regions real test images never occupy. This
species-specific gap is the structural cause of the downstream
behaviour in \cref{sec:downstream}.

\subsection{LLM-as-Judge Results}
\label{sec:llm-results}

\Cref{tab:llm} reports the two-stage LLM judge on the full
$1{,}500$-image generated pool. Blind-ID rates are above $96\%$ for
\emph{B.~sandersoni} and \emph{B.~flavidus} but only $76\%$ for
\emph{B.~ashtoni}; misidentified \emph{ashtoni} cases concentrate among
yellow-thorax congeners (\emph{B.~impatiens}, \emph{B.~griseocollis}),
confirming that both generation \emph{and} judgment inherit the
worker-bumblebee prior. Strict pass rates are strongly
species-dependent: $44.4\%$, $91.2\%$, and $57.6\%$.

\begin{table}[t]
  \caption{LLM-as-judge on $1{,}500$ generated images. Strict pass
  requires blind-ID match, species-level diagnosis, and mean morph
  $\geq 4.0$.}
  \label{tab:llm}
  \centering
  \setlength{\tabcolsep}{6pt}
  \begin{tabular}{@{}lrrrr@{}}
    \toprule
    Species & Blind-ID & Morph & Diag.=sp. & Strict \\
    \midrule
    \emph{B.~ashtoni}    & 76.0\% & 3.82 & 59.6\% & 44.4\% \\
    \emph{B.~flavidus}   & 96.0\% & 4.06 & 61.2\% & 57.6\% \\
    \emph{B.~sandersoni} & 96.4\% & 4.37 & 91.2\% & 91.2\% \\
    \bottomrule
  \end{tabular}
\end{table}

A per-feature morphology heatmap (supplementary) shows $14$ of the $15$
species-by-feature cells exceed $4.0$; the lone outlier is
\emph{B.~ashtoni} thoracic coloration ($2.98$), and \emph{wrong
colouration} is the dominant failure flag ($27.1\%$ of images,
concentrated on \emph{ashtoni}). Structural flags (impossible geometry,
missing/extra limbs, artefacts) are essentially absent, so structured
prompting removes gross anatomical errors but not species-specific
colour errors.

\subsection{Expert Calibration}
\label{sec:expert-results}

To ground these automated signals in biological validity, the
entomologist scored the $150$-image stratified sample under the same
two-stage rubric. Only $50/150$ images pass the strict expert
rule (against $128/150$ for the lenient rule), and the rate varies
sharply by species: $27/50$ for \emph{B.~sandersoni} and $17/50$ for
\emph{B.~ashtoni}, but just $6/50$ for \emph{B.~flavidus}---the clearest
case in which general visual plausibility diverges from species-faithful
morphology. Rejections are dominated by ``does not match target
species,'' followed by wrong colouration and structural issues, and
free-text notes locate the ``species other'' cases in fine-grained
anatomy: face shape, antennae, mouthparts, leg morphology, and body
proportions.

Measured against this reference, the LLM judge is only weakly
calibrated. Its strict rule passes $78$ images to the expert's $50$, but
the two accept the same $31$ (LLM precision $0.40$, recall $0.62$);
overall agreement is $84/150$ ($56\%$), barely seven points above what
marginal independence predicts, and as a continuous ranker the LLM
morphology score reaches only AUC $0.56$. The disagreement is structured
rather than random: the judge is over-strict on \emph{B.~ashtoni}'s dark
thorax, which the expert accepts, yet misses subtle errors in
\emph{B.~flavidus} colouration, head shape, and appendages
(\cref{fig:expert}).

\begin{figure}[t]
  \centering
  \includegraphics[width=0.58\linewidth]{expert/llm_vs_expert_strict_2x2.png}
  \caption{LLM strict versus expert strict on the $150$ annotated
  images.}
  \label{fig:expert}
\end{figure}

A logistic probe on frozen BioCLIP, supervised by the same $150$ labels,
recovers a far stronger signal. Among four input configurations
(BioCLIP alone; LLM-derived features alone; their concatenation; and the
concatenation with a species one-hot), BioCLIP alone attains the highest
LOOCV AUC ($0.792$) under the strict rule and the LLM features add
nothing (\cref{fig:probe}\,(a)), so the deployed D6 probe uses BioCLIP
alone with per-species F1-maximising thresholds. \emph{B.~sandersoni}
separates cleanest, while \emph{B.~ashtoni} and \emph{B.~flavidus} keep
the overlap expected from their lower strict-pass rates. The practical
payoff shows up at fixed budget (\cref{fig:probe}\,(b)): counting
expert-strict positives inside each filter's deployed $200$-image
selection, the D4 LLM rule retains $5/12/1$ images for
\emph{ashtoni}/\emph{sandersoni}/\emph{flavidus} and the D5 centroid
$9/13/1$, whereas D6 retains $16/20/6$---roughly $3\times$ richer in
expert-validated images under the same cap.

\begin{figure}[t]
  \centering
  \begin{subfigure}{0.49\linewidth}
    \centering
    \includegraphics[width=\linewidth]{expert/probe_roc_loocv.png}
    \caption{Per-species LOOCV ROC (overall AUC $0.79$); red markers
    are the F1-maximising deployment thresholds.}
  \end{subfigure}
  \hfill
  \begin{subfigure}{0.49\linewidth}
    \centering
    \includegraphics[width=\linewidth]{expert/expert_coverage_of_selected200.png}
    \caption{Expert-strict composition of each filter's $200$-image
    selection at fixed budget.}
  \end{subfigure}
  \caption{Expert-calibrated probe (D6): probe ROC (a) and, under
  volume parity, the expert-strict composition of each filter's
  $200$-image selection (b).}
  \label{fig:probe}
\end{figure}

A single image makes the gap concrete (\cref{fig:disagreement}). One
\emph{B.~flavidus} synthetic passes the LLM, centroid, and probe filters
(mean morph $4.20$, matching blind ID, centroid $0.785\!\ge\!0.682$) yet
the expert rejects it for a missing antenna; a second fails every
automated filter (the LLM calls it \emph{B.~ternarius}, centroid
$0.604\!<\!0.682$) yet the expert accepts it on a correct blind ID
because the viewpoint hides the head. Language-mediated reasoning,
feature-space proximity, and expert validity disagree image by image,
not only in aggregate.

\begin{figure}[t]
  \centering
  \includegraphics[width=0.72\linewidth]{expert/flavidus llm vs automated.png}
  \caption{\emph{B.~flavidus} disagreement pair: one synthetic passes
  all automated filters but fails the expert (missing antenna), the
  other fails every filter but the expert accepts it (viewpoint hides
  the head).}
  \label{fig:disagreement}
\end{figure}

\subsection{Downstream Augmentation Comparison}
\label{sec:downstream}

\Cref{tab:downstream} reports ResNet-50 macro-F1 across the six dataset
variants and three protocols. Three patterns dominate: (i)~augmentation
effects concentrate in the rare tier, with moderate and common species
barely changing; (ii)~the protocols give different aggregate rankings
(D6 in the single run, baseline under multi-seed, copy-paste under
five-fold CV), so none should be over-interpreted; and (iii)~within the
generated-filter family D6 is the most consistent, best or tied-best
across all three views.

\begin{table}[t]
  \caption{Macro-F1 across six dataset variants and three protocols.
  Rare-F1 is reported as \emph{ash}/\emph{sand}/\emph{flav}.
  Per-protocol best macro-F1 is in \textbf{bold}; multi-seed and
  five-fold values are $\mathrm{mean}\pm\mathrm{sd}$.}
  \label{tab:downstream}
  \centering
  \resizebox{\linewidth}{!}{%
  \setlength{\tabcolsep}{4pt}
  \begin{tabular}{@{}l ccc ccc ccc@{}}
    \toprule
    & \multicolumn{3}{c}{\textbf{Single split}}
    & \multicolumn{3}{c}{\textbf{Multi-seed (5)}}
    & \multicolumn{3}{c}{\textbf{Five-fold CV}} \\
    \cmidrule(lr){2-4}\cmidrule(lr){5-7}\cmidrule(lr){8-10}
    \textbf{Variant}
      & F1 & 95\% CI & Rare F1
      & F1 & Rare F1 & Acc
      & F1 & Rare F1 & Acc \\
    \midrule
    D1 Baseline
      & 0.815 & [.774,.845] & .50/.59/.62
      & \textbf{0.839$\pm$0.006} & .61/.62/.76 & 0.890
      & 0.832$\pm$0.013 & .62/.47/.75 & 0.892 \\
    D2 CNP
      & 0.829 & [.787,.860] & .55/.63/.72
      & 0.822$\pm$0.014 & .58/.48/.72 & 0.886
      & \textbf{0.837$\pm$0.013} & .69/.43/.81 & 0.891 \\
    D3 Unfiltered
      & 0.823 & [.783,.854] & .50/.53/.70
      & 0.828$\pm$0.009 & .61/.49/.67 & 0.891
      & 0.820$\pm$0.024 & .58/.33/.76 & 0.889 \\
    D4 LLM
      & 0.834 & [.791,.864] & .60/.59/.71
      & 0.831$\pm$0.008 & .61/.53/.71 & 0.890
      & 0.821$\pm$0.019 & .58/.40/.74 & 0.889 \\
    D5 Centroid
      & 0.824 & [.782,.855] & .55/.50/.73
      & 0.827$\pm$0.002 & .62/.49/.70 & 0.889
      & 0.824$\pm$0.020 & .53/.38/.76 & 0.893 \\
    D6 Expert-probe
      & \textbf{0.840} & [.801,.869] & .67/.56/.75
      & 0.834$\pm$0.008 & .64/.55/.70 & 0.890
      & 0.824$\pm$0.016 & .60/.38/.75 & 0.892 \\
    \bottomrule
  \end{tabular}}
\end{table}

\subsubsection{Statistical Tests.}
Paired $t$-tests on fold-level macro-F1 yield only \emph{negative}
significant aggregate effects: LLM filtering (D4) is significantly worse
than both the baseline ($\Delta=-0.011$, $p=0.04$) and copy-paste
($\Delta=-0.016$, $p=0.03$), and no generated-filter variant
significantly beats the unfiltered pool (full pairwise table in the
supplementary). The sole significant
\emph{positive} effect is per-species---copy-paste (D2) improves
\emph{B.~flavidus} F1 over baseline---making copy-paste the most
statistically supported augmentation in this study.

\subsubsection{Volume Ablation: Quantity Is Not the Bottleneck.}
Increasing the synthetic budget from $+50$ to $+300$ images per rare
species under D3--D6 produces no monotonic improvement and all bootstrap
CIs overlap the baseline: adding more images of the same quality does
not help, so the binding constraint is selection, not scale.

\subsubsection{Subtractive Ablation Locates Harm.}
Removing all synthetics for one rare species at a time and remeasuring
own-species F1 (\cref{fig:subtract}) reveals where each filter's pool
was harmful or helpful: a negative bar means removal \emph{decreased}
F1, so the removed images were collectively helpful. D3 unfiltered
synthetics harm both \emph{B.~ashtoni} and \emph{B.~flavidus}; D4 LLM
helps \emph{B.~ashtoni} but harms \emph{B.~sandersoni}; D5 centroid
leaves harmful pools for two of three targets. D6 is the only variant
for which all three own-species cells are non-positive, indicating its
selections are net helpful for every rare species.

\begin{figure}[t]
  \centering
  \includegraphics[width=0.6\linewidth]{ablation/subset_ablation_recovery_bars.png}
  \caption{Own-species F1 recovery under single-species subtractive
  ablation. Bars below zero (green) mark synthetics that were
  collectively helpful, above zero (red) harmful; D6 is the only
  variant with all three rare-species cells non-positive.}
  \label{fig:subtract}
\end{figure}

\subsubsection{Per-Image Flips.}
Collapsing multi-seed predictions by majority vote and comparing against
D1 categorises each test image as stable-correct, stable-wrong,
improved, or harmed. Most rare-species \emph{changes} are harms, not
improvements; D6 is the only variant with zero harmed images for both
\emph{B.~ashtoni} and \emph{B.~sandersoni} while rescuing one
\emph{B.~flavidus} image. BioCLIP nearest-neighbour retrieval from each
harmed image shows the nearest training synthetic is often far in
absolute terms and sometimes a confuser species, making the
\cref{sec:gap} feature gap visible at image level.


% ===============================================================
\section{Discussion and Conclusion}
\label{sec:discussion}

\subsubsection{Fidelity, Not Volume, Is the Binding Constraint.}
Across protocols, copy-paste is the only augmentation with a
statistically supported \emph{per-species} F1 improvement over the
baseline (\emph{B.~flavidus}, five-fold CV); its macro-F1 gain is
directional but not separable ($p=0.16$), and no generated-filter
variant beats the unfiltered pool by a separable margin. Volume
ablation rules out the simplest explanation: $+50$ to $+300$ generated
images per rare species produce no monotonic improvement. What binds is
whether images preserve expert-valid, classifier-relevant morphology,
not how many exist. This gap is species-specific---each rare target's
synthetic cluster occupies its own displaced region of BioCLIP space
with a distinct failure profile---so one-size-fits-all augmentation is
insufficient for fine-grained biodiversity.

\subsubsection{Synthetics Reshape Rare-Tier Boundaries.}
A single-seed additive ablation (supplementary) shows synthetic images
do not merely add examples to one class but reshape decision boundaries
among confusable rare classes: $+200$ \emph{B.~sandersoni} synthetics
help mainly by improving \emph{B.~ashtoni}, while \emph{B.~flavidus}
synthetics help their own class but harm the others. Augmentation should
thus be evaluated across the rare tier, not per-class.

\subsubsection{Automated Quality Signals Are Not Expert Validity.}
The central calibration finding is that the LLM judge and the expert
diverge substantially and \emph{systematically}---through viewpoint
occlusion, species-specific complexity, LLM blind spots, and centroid
proximity conflating pose, crop, and lighting with species identity.
Supervising a probe on the expert labels recovers much of what these
unsupervised signals miss, and its three ingredients (a stratified
expert sample, a probe on a frozen domain-aligned embedding, and
per-class F1-maximising thresholds) are not bumblebee-specific: they
could plausibly transfer to other fine-grained long-tailed categories,
though we leave that test to future work.

\subsubsection{Deployment Implications.}
For monitoring systems such as camera-trap networks and city-scale
insect-monitoring platforms~\cite{venverloo_towards_2024}, a
rare-species miss can become a false ecological absence in the record. The conservative recipe is to
prefer copy-paste where source images exist, gate generative
augmentation on expert-calibrated filters, and route rare-species
predictions through a human-in-the-loop channel---tracking the
operational \emph{false-absence} rate against ground-truth surveys
rather than held-out F1 alone.

\subsubsection{Limitations.}
Three constraints bound our conclusions. First, the calibration target
is a single expert: the $150$ labels define a reference rubric rather
than consensus ground truth, so we cannot estimate inter-annotator
agreement or separate idiosyncratic from systematic judgments, and the
probe inherits whatever biases the rubric encodes. Second, the
downstream comparison is underpowered at the rare-tier scale
($n{=}6$, $10$, $36$ test images; \cref{sec:experiments}): single
prediction flips dominate per-species F1, and five-fold CV leaves only
four degrees of freedom for paired tests, so our downstream claims are
deliberately conservative. Third, the study fixes one classifier
backbone (ResNet-50), one generator (\texttt{gpt-image-1.5}), one taxon
(Massachusetts \emph{Bombus}), and one frozen embedding (BioCLIP, a
biologically informed but imperfect proxy for the classifier's own
representation); none of these is guaranteed to generalise, and BioCLIP
distances may not coincide with the boundaries the classifier learns.

\subsubsection{Future Work.}
The strongest near-term recipe is copy-paste where source images exist,
complemented by selective generation filtered through expert-calibrated
visual representations. Natural extensions are scaling the
expert-annotation set across multiple entomologists, replicating the
D1--D6 comparison under biology-aware
backbones~\cite{stevens_bioclip_2024, oquab_dinov2_2024} and on
classifier-internal features, fine-tuning open-source diffusion models
on entomological imagery (\eg via LoRA) to close the synthetic--real gap
at the source, and evaluating the pipeline against held-out ecological
surveys using a false-absence rate rather than held-out F1.


% ---- Bibliography ----
% BibTeX users should specify bibliography style 'splncs04'.
\bibliographystyle{splncs04}
\bibliography{main}

\end{document}
